\begin{lemma}
For every positive integer $k$ there is $D_0$ such that the following holds
for all $D\ge D_0$.  Let $\mathcal F\in H(k,k,D)$, and let
$f\in E(\mathcal F)$ maximize $b_{L(\mathcal H),kD}(e)$ over all
$\mathcal H\in H(k,k,D)$ whose vertex and edge types lie in the same
ambient universe levels as those of $\mathcal F$, and all
$e\in E(\mathcal H)$.  Then
\[
P_{L(\mathcal F)}(f)\ge \frac{k-1}{2}D^2-k^2D .
\]
\end{lemma}
\begin{proof}
Choose $D_0$ large enough that $D\ge D_0$ implies $D\ge 1$ and that
\noderef{KUniformKSimpleExtremalSaturation} applies.  Fix $D\ge D_0$.
We first dispose of the case $k=1$.  By \noderef{IndependentPairCount},
$P_{L(\mathcal F)}(f)$ is the cardinality of a set of independent
two-subsets of the neighborhood of $f$ in the line graph.  Hence
$P_{L(\mathcal F)}(f)\ge 0$.  Since $D\ge 1$, also $D\ge0$, and when
$k=1$ the required right hand side is
\[
\frac{k-1}{2}D^2-k^2D=-D\le 0 .
\]
Thus $P_{L(\mathcal F)}(f)\ge0\ge -D$, proving the claim for $k=1$.

It remains to treat $k\ge2$.
Let $f=\{v_1,\ldots,v_k\}$ and let $X=X(f,\mathcal F)$ be the set of
vertices outside $f$ that occur in an edge meeting $f$.  Let $G$ be the
subgraph of \noderef{LineGraphOfHypergraph} induced on
$N_{L(\mathcal F)}(f)$.  By
\noderef{KUniformKSimpleExtremalSaturation}, every vertex $v_i$ of $f$ has
degree $D$, and $f$ has exactly $k(D-1)$ line-graph neighbors.  Hence the
total number of unordered pairs in $N_{L(\mathcal F)}(f)$ is
\[
\binom{kD-k}{2}.
\]
By \noderef{IndependentPairCount}, $P_{L(\mathcal F)}(f)$ is the number of
nonedges among these pairs.  Therefore
\[
P_{L(\mathcal F)}(f)=\binom{kD-k}{2}-|E(G)|. \tag{1}
\]

We next bound $|E(G)|$.  For each $i$, let
$C_i$ be the set of edges other than $f$ that contain $v_i$.  The saturation
node says that the $C_i$ are disjoint and each has size $D-1$: disjointness
follows because the total incidence count from the $k$ vertices of $f$ is
already the number of distinct neighbors, so no neighbor can contain two
vertices of $f$; then the maximum-degree bound forces each class to have
exactly $D-1$ members.  Inside each $C_i$ all pairs are adjacent in the line
graph, contributing
\[
k\binom{D-1}{2}. \tag{2}
\]

For cross-edges between $C_i$ and $C_{i'}$, put
$a_{i,x}=\deg_{\mathcal F}(v_i,x)$ for $x\in X$.  A cross pair
$(e_1,e_2)\in C_i\times C_{i'}$ can be adjacent only by sharing a vertex of
$X$, because $e_1$ contains $v_i$ but not $v_{i'}$ and $e_2$ contains
$v_{i'}$ but not $v_i$.  Summing
$a_{i,x}a_{i',x}$ over $x\in X$ counts such a cross pair once for each
outside common vertex.  Thus it is an upper bound for the number of
line-graph edges between $C_i$ and $C_{i'}$.  Summing over $i<i'$ gives
\[
|E(G)|\le
k\binom{D-1}{2}
+
\sum_{x\in X}\sum_{1\le i<i'\le k}a_{i,x}a_{i',x}. \tag{3}
\]

The table $A=(a_{i,x})$ satisfies the hypotheses of
\noderef{TablePairFillingBound} with $C=D$ and $T=k(k-1)D$.  Indeed, each
entry is a nonnegative degree count by \noderef{HypergraphDegree}.  The
already established assumptions for this remaining case give the two
numerical side conditions in that node: $k\ge2$ by the case split,
$C=D>0$ because $D\ge1$, and
\[
T=k(k-1)D>0
\]
because each of the three integer factors $k$, $k-1$, and $D$ is positive.
For a
fixed outside vertex $x$, the column sum
$\sum_i a_{i,x}$ counts edges containing $x$ and one vertex of $f$; by the
disjointness just proved, no such edge is counted twice, so this column sum
is at most the degree of $x$, and hence at most $D$ by
\noderef{MaxDegreeAtMost}.  For a fixed row $i$, the sum
$\sum_x a_{i,x}$ counts outside incidences in the $D-1$ edges of
$C_i$.  Each such edge is $k$-uniform by \noderef{UniformHypergraph} and
contains $v_i$ plus exactly $k-1$ vertices outside $f$, so the row sum is
$(k-1)(D-1)$.  Therefore the whole table sum is
$k(k-1)(D-1)\le k(k-1)D$.

Applying \noderef{TablePairFillingBound} with this $C$ and $T$ gives
\[
\sum_{x\in X}\sum_{i<i'}a_{i,x}a_{i',x}
\le
\frac{k-1}{2k}\,k(k-1)D^2
=\frac{(k-1)^2}{2}D^2. \tag{4}
\]
Combining (1), (3), and (4), we obtain
\[
\begin{aligned}
P_{L(\mathcal F)}(f)
&\ge
\binom{kD-k}{2}
-k\binom{D-1}{2}
-\frac{(k-1)^2}{2}D^2\\
&=
\frac{k-1}{2}D^2-(k^2-k)D+\frac{k^2-k}{2}\\
&\ge \frac{k-1}{2}D^2-k^2D,
\end{aligned}
\]
which is the asserted lower bound.
\end{proof}
